\section{讨论}

\subsection{神经上下文的局部性}

在两个数据集上，少量邻近词窗已经能够带来较大的候选排序改善，继续扩大上下文范围并未持续提高主指标。
在固定 $50$ 词候选集上，局部范围对宏平均 Top-$10$ 更有利，而完整训练词库下的 Top-$1$ 不呈现一致优势，
说明不同评价尺度对上下文范围的响应存在差异。
长窗单词模型未达到当前局部跨词模型的检索水平，但这并不构成对跨词神经信息来源的单独识别。
本文更关注有限上下文在既定输入与模型协议下的实用表现，以及随后可实现的质量—等待折中。

\subsection{词级与序列级结果的差异}
\label{sec:discussion_word_sentence}

局部窗口在词级宏平均 Top-$10$ 上占优，而在法语 MEG 的最终文本质量比较中，完整片段策略 F 获得更高的 BLEU-1。
宏平均检索对词类型等权，文本指标则同时受到词频、候选构造和序列搜索规则影响，两者不对应相同的优化目标。
法语的 F-local 与 C 的最终文本质量接近，使用完整上下文模型的策略 F 进一步提高 BLEU-1；
F 和 F-local 来自不同上下文范围下独立训练的模型。
中文 EEG 各策略的文本差异较小。
上述实验用于比较不同输入等待与搜索协议下的预测表现，不将数据集之间的差异归因为单一神经或语言机制。

\subsection{滚动修正与输出稳定性}

多数回修仍发生在错误词之间，输出稳定性需要与最终文本质量分别评价。
后续工作可根据候选得分边际或置信度触发回修，并采用重翻译研究中的归一化擦除指标评价显示稳定性\cite{arivazhagan2020retranslation}。

\subsection{局限}

本文仍有若干适用边界。
首先，实验基于已知词级起始时刻、片段边界和目标词数的离线顺序重放，
输入等待仅反映神经窗口可用时刻，不包含模型推断与文本搜索耗时。
其次，相邻词窗高度重叠，模型性能可能受到词间时间间隔及其他输入结构线索的影响；
本文不对这些线索与词相关脑活动的贡献作定量分离。
再次，测试受试者出现在训练阶段，因此当前结果反映已见受试者上的新材料表现，文本绝对词位恢复率也仍然有限。
最后，两个数据集均来自《小王子》，语言模型可能在预训练阶段接触相关文本；
本文不将最终文本得分解释为神经信息贡献的独立估计。
